% Input by report.tex. All mathematical abbreviations in this section are local.
\section{A concrete equation in seventeen unrestricted integers}
\label{sec:integer17}

This section gives a fixed polynomial $P_{17}$ in seventeen unrestricted
integer unknowns whose absence of zeros is equivalent, provably in ZFC, to
$\Con(\mathrm{ZFC})$. It uses the complete original T4 coefficient tuple of
package 18, \texttt{18-jones9}, at revision \pin\ \cite{jones9}.
All source coefficients and abbreviations below are local to this section.
In particular, this construction has its own coefficient binding, distinct
from the $943$-tile binding used in the preceding sections.

The result improves the delivered concrete $26$-integer count in the
inspected direct compiler family, at the cost of an enormous degree.
It establishes neither global minimality nor a universal-variable record:
Sun's ten-integer representation theorem already gives a smaller abstract
general bound \cite{sun2026}. The contribution here is the complete fixed
specialization, its explicit circuit, and the additional proofs needed to
use signed coordinates. The source Jones ratio theorem, quartic/compiler
theorem, and K215 calibration are inherited ordinary mathematical theorems.
The new arguments below are proved in full; no new proof-assistant
certificate is claimed.

\subsection{The fixed coefficients and exact source interface}
\label{i17:source}

Start with eight base coordinates
\[
 (\varepsilon,g_0,h_0,s_0,w_0,\phi,i_0,j_0)\in\Z^8
\]
and abbreviate
\[
 g=g_0+1,\quad h=h_0+1,\quad s=s_0+1,\quad
 w=w_0+1,\quad i=i_0+1,\quad j=j_0+1.
\]
The old ninth coordinate of package 18 and its finalizer are replaced
below. The full source expressions through the ratio conditions are
reproduced here so that the signed-domain argument has an explicit
interface.

\paragraph{Finite coefficient recipe.}
Put $\nu=32$ and $L=5^{33}$. Shift the $32$ positive coordinates of the
pinned quartic by $1$, and include its unused dummy input as coordinate
zero. The dummy is fixed to $1$ and is not an unknown of $P_{17}$.
Let $p_\alpha$ be the resulting coefficient indexed by
$\alpha\in\N^{33}$, with
$|\alpha|=\sum_{r=0}^{32}\alpha_r\le4$, and put
\[
 c_\alpha=(4-|\alpha|)!\left(\prod_{r=0}^{32}\alpha_r!\right)p_\alpha.
\]
The four compiler ports are bound to the complete original T4 tuple in
\path{signed17/upstream/COEFFICIENT_BINDINGS.json}.
The coefficient records and the positive-coefficient bounds $\rho_j$ for
the shifted source residuals are specified by
\path{signed17/upstream/nine/generated/constants.json} and
\path{signed17/upstream/nine/generated/index_bound_s.json}.
Define the following fixed integers and fixed polynomial:
\begin{align*}
 s_\star&=2+24\sum_j\rho_j^2,
 &z_\star&=2^{s_\star},\\
 e_\star(T)&=
 \sum_{\substack{\alpha\in\N^{33}\\|\alpha|\le4}}
 (z_\star+c_\alpha)T^{\,L-\sum_{r=0}^{32}\alpha_r5^r},
 &y_\star&=e_\star(2z_\star).
\end{align*}
The sum includes the unlisted zero coefficients. All exponents are fixed
nonnegative integers. The source bound proves
\[
 z_\star>|c_\alpha|+1\quad\text{for every }\alpha,
 \qquad y_\star\ge2z_\star.
\]
Thus the coefficient recipe is fixed independently of the quantified
unknowns and of the consistency of ZFC.

\paragraph{Source polynomial abbreviations.}
Define
\begin{align*}
 b&=y_\star+1+\varepsilon,
 &\mathcal B&=b^5,
 &q&=\mathcal B^L,\\
 \ell&=\sum_{r=1}^{\nu}\mathcal B^{5^r},
 &\lambda&=\sum_{r=0}^{4L-1}\mathcal B^r,
 &e&=e_\star(\mathcal B),\\
 \theta&=\mathcal B-2z_\star,
 &N&=q^{16},
 &G&=32b\mathcal B^{5^\nu}.
\end{align*}
To display the packed source expression without a long line, write
\begin{align*}
 R_1={}&g+eq^3+\ell q^5\\
 &+\bigl[2(e-z_\star\lambda)(1+\mathcal B+g)^4
       +\lambda\mathcal B(1+q^4)\bigr]q^7,\\
 R_2={}&q^3-b\ell+\ell+\theta\lambda q^3+(\mathcal B-2)q^8,\\
 R={}&R_1(N^2-N)+R_2(N^2-1).
\end{align*}
The remaining abbreviations are
\begin{align*}
 U&=N^2w,&Y&=N^2s,&M&=RY,\\
 P&=2M^2U,&K&=R+1+h(P-1),&A&=M(U+1),\\
 B_1&=2R+1,&W_1&=bw,&C&=B_1+W_1+\phi,\\
 J&=4A-5,&Q&=3W_1C-2(W_1^2-1),
\end{align*}
and
\begin{align*}
 D&=(A^2-1)C^2+1,&E&=iJDC^2,\\
 F&=(A^2-1)E^2+1,&G_1&=A+F(CD+1-A),\\
 H&=B_1+jC,&I&=(G_1^2-1)H^2+1,\\
 a_0&=(P^2-1)K^2+1,&a_1&=DFI,\\
 D_0&=JF,&T_0&=QF+J(H-C).
\end{align*}
Finally, set
\[
 m=K^2G-4(C-KY)^2G-8gK^2.
\]
Let $\operatorname{Sq}(\Z)=\{r^2:r\in\Z\}$.
The imported Jones ratio equivalence states that there exist eight
\emph{nonnegative} base coordinates satisfying
\begin{equation}
 a_0,a_1\in\operatorname{Sq}(\Z),\qquad
 D_0\mid T_0,\qquad m>0
 \label{i17:JR}
\end{equation}
if and only if the pinned $32$-variable quartic has a positive solution.
The imported compiler and K215 calibration identify that assertion with
existence of a finite ZFC refutation. These are the precise inherited
theorem boundaries used below. The signed-coordinate argument depends
on the displayed source structure, not merely on existence of an
unspecified nine-variable representation.

\subsection{Two nonnegativity guards are redundant}

Only six base coordinates need explicit nonnegativity guards:
\[
 \mathbf x=(x_1,\ldots,x_6)
          =(\varepsilon,g_0,s_0,w_0,\phi,j_0).
\]

\begin{lemma}[The source denominator never vanishes]
\label{i17:denominator}
For every integer assignment to the eight base coordinates, $D_0\ne0$.
\end{lemma}
\begin{proof}
The integer $J=4A-5$ cannot vanish. If $F=0$, then
$(A^2-1)E^2=-1$, which forces $A=0$ and $E=\pm1$.
But $A=0$ gives $J=-5$, so $E=iJDC^2$ is divisible by $5$, a contradiction.
\end{proof}

\begin{lemma}[A signed coordinate]
\label{i17:signed-i}
If $i\ne0$, replacing $i$ by $|i|$ leaves
$a_0,a_1,D_0,T_0,m$ unchanged.
\end{lemma}
\begin{proof}
The coordinate $i$ enters through $E$, and these five expressions use
$E$ only through $E^2$ in $F$. The replacement therefore preserves all
five values. The new base coordinate $i_0=|i|-1$ is nonnegative.
\end{proof}

\begin{lemma}[The guard on $h_0$ follows from the source conditions]
\label{i17:h-guard}
Suppose that the six coordinates $\mathbf x$ are nonnegative,
$a_0\in\operatorname{Sq}(\Z)$, and $m>0$. Then $h\ge1$, so $h_0\ge0$.
\end{lemma}
\begin{proof}
The coefficient recipe gives $b\ge2$, $\mathcal B=b^5\ge32$,
$e\ge0$, and $\mathcal B>2z_\star$. Put $t=5^\nu$.
In particular, $G>0$ and $g\ge1$. The positive margin first implies
\begin{equation}
 K\ne0,\qquad g<G/8=4b\mathcal B^t,\qquad
 4(C-KY)^2<K^2.
 \label{i17:margin-bounds}
\end{equation}
Indeed, $K=0$ would make $m\le0$, while
$m\le(G-8g)K^2$ gives the bound on $g$.
Deleting the strictly positive term $8gK^2$ from the margin inequality
and dividing by $G$ gives the last inequality.

Since $b\le\mathcal B$ and $\mathcal B>5$, we have
\begin{equation}
 \begin{aligned}
 1+\mathcal B+g
 &<1+\mathcal B+4\mathcal B^{t+1}
 \le5\mathcal B^{t+1}\\
 &<\mathcal B^{t+2}\le\mathcal B^{5t}=q.
 \end{aligned}
 \label{i17:g-bound}
\end{equation}
The distinct exponents in $\ell$ lie between $0$ and $t$, so
\[
 \ell\le\sum_{r=0}^t\mathcal B^r<\mathcal B^{t+1},
 \qquad b\ell<\mathcal B^{t+2}\le q<q^3.
\]
Writing $V_0=1+\mathcal B+g$, the possibly negative part of $R_1$ is
dominated by its positive companion:
\[
 \lambda\mathcal Bq^4-2z_\star\lambda V_0^4>0.
\]
This follows from \eqref{i17:g-bound} and $\mathcal B>2z_\star$.
After separating $2(e-z_\star\lambda)V_0^4$ into its two terms,
all remaining terms of $R_1$ are nonnegative, and $g>0$.
Also $R_2>0$, since $q^3>b\ell$, $\theta>0$, and $\mathcal B>2$.
Consequently
\[
 R\ge N^2-N\ge2.
\]
This proves positivity of $R$ from the margin-derived bound, without
assuming any sign for $h$.

Now $Y,U\ge1$, $C=2R+1+bw+\phi>0$, and
$P=2R^2Y^2U\ge2R^2>R+2$.
If $K<0$, then
$|C-KY|=C+|K|Y>|K|$, contrary to \eqref{i17:margin-bounds}.
The case $K=0$ has already been excluded, so $K>0$.

Suppose $h\le0$. Then $1\le K\le R+1<P$.
If $K=1$, the identity $K=R+1+h(P-1)$ gives $P-1\mid R$,
contrary to $0<R<P-1$.
If $2\le K<P$, then
\[
 (PK-1)^2<(P^2-1)K^2+1<(PK)^2.
\]
The difference on the left is $K(2P-K)>0$, and the difference on
the right is $K^2-1>0$. Thus $a_0$ lies strictly between consecutive
squares, a contradiction. Hence $h\ge1$.
\end{proof}

It follows that the existential predicate \eqref{i17:JR} is unchanged
if $h_0$ and $i_0$ range over all integers, provided $i=i_0+1$ is
nonzero and the six coordinates $\mathbf x$ remain nonnegative.

\subsection{Seven inequalities from one square condition}

Introduce seven integer unknowns $z_1,\ldots,z_7$. To avoid confusing
the following factors with the source quantities $B_1$ and $C$, use
the symbols $\eta_j$ and $\Gamma_j$.
Define
\begin{equation}
 \begin{aligned}
 \Gamma_0&=1,\\
 d_j&=3x_j+2 &&(1\le j\le6),\\
 d_7&=3m-1,\\
 \eta_j&=1+4d_j\Gamma_{j-1}^2z_j^2,
 &\Gamma_j&=\Gamma_{j-1}\eta_j\qquad(1\le j\le7).
 \end{aligned}
 \label{i17:stack}
\end{equation}

\begin{lemma}[A coprime Pell stack]
\label{i17:pell-stack}
Subject to $z_1\cdots z_7\ne0$, the condition
$\Gamma_7\in\operatorname{Sq}(\Z)$ implies
$x_1,\ldots,x_6\ge0$ and $m>0$.
Conversely, if these seven inequalities hold, there are nonzero
$z_1,\ldots,z_7\in\Z$ for which $\Gamma_7$ is a square.
\end{lemma}
\begin{proof}
Every $\eta_j$ is $1$ modulo $4$, hence nonzero, and
$\gcd(\eta_j,\Gamma_{j-1})=1$. Thus the seven factors $\eta_j$ are
pairwise coprime. If their product is a square, unique factorization
implies that each $|\eta_j|$ is a square. The negative of an integer
square is $0$ or $3$ modulo $4$, so $\eta_j\equiv1\pmod4$ forces
each $\eta_j$ to be a positive square.

Both $\Gamma_{j-1}$ and $z_j$ are nonzero. If $d_j<0$, then
$\eta_j\le1-4=-3$, a contradiction. Hence $d_j\ge0$.
Since $d_j\equiv2\pmod3$, equality is impossible. Therefore
$d_j>0$, which is exactly $x_j\ge0$ for $1\le j\le6$ and
$m\ge1$ for $j=7$.

Conversely, every $d_j$ is then positive and congruent to $2$ modulo
$3$, so it is not a square. The Pell equation
$u_j^2-d_jv_j^2=1$ has a nontrivial positive solution.
It has one whose $v_j$ is divisible by any prescribed positive integer
$M_0$: powers of the associated unit matrix
\[
 \begin{pmatrix}u_j&d_jv_j\\v_j&u_j\end{pmatrix}
 \in\mathrm{SL}_2(\Z)
\]
eventually become the identity modulo $M_0$, because
$\mathrm{SL}_2(\Z/M_0\Z)$ is finite. A positive power remains
nontrivial and has positive lower-left entry. Set
$M_0=2|\Gamma_{j-1}|$ and choose such a solution. Then
\[
 z_j=\frac{v_j}{2\Gamma_{j-1}}\in\Z\setminus\{0\},
 \qquad \eta_j=u_j^2.
\]
Induction gives seven positive-square factors and a square product.
\end{proof}

The nonzero conditions on the seven witnesses will be imposed together
by one divisibility relation. In particular, no zero witness is being
excluded without an algebraic constraint.

\subsection{Divisibility and nonvanishing through a degree-five polynomial}

Let
\[
 \Pi(t)=(t^3-t-1)(t^2+23).
\]

\begin{lemma}[Intersectivity]
\label{i17:intersective}
The polynomial $\Pi$ has no rational root, but for every nonzero
integer $n$ there is an integer $t$ such that $n\mid\Pi(t)$.
\end{lemma}
\begin{proof}
The monic cubic $f(t)=t^3-t-1$ has no rational root: the only candidates
are $1$ and $-1$, and neither is a root. The quadratic has no real root.

It remains to construct a root modulo every prime power. The cubic has
discriminant $-23$. For an odd prime $p\ne23$, any root of the cubic
modulo $p$ is simple and therefore lifts to all powers of $p$.
If it has no root modulo $p$, it is irreducible of degree three.
Its roots lie in $\mathbb F_{p^3}$, and Frobenius permutes them by a
three-cycle. The product of their pairwise differences is fixed by this
even permutation, so lies in $\mathbb F_p$; its square is $-23$.
It is nonzero. Thus the quadratic $t^2+23$ has a simple root modulo
$p$, which also lifts to every power of $p$.
For $p=23$, the cubic has the simple root $3$, because
$f(3)=23$ and $f'(3)=26$.

For $p=2$, use $-23\equiv1\pmod8$.
If an odd $t$ satisfies $t^2\equiv-23\pmod{2^e}$ with $e\ge3$,
then one of $t$ and $t+2^{e-1}$ satisfies the same congruence modulo
$2^{e+1}$: their squares differ by $2^et$ modulo $2^{e+1}$.
The initial cases modulo $2$, $4$, and $8$ are immediate. This supplies
roots modulo every power of $2$. The Chinese remainder theorem now
gives a root modulo $|n|$.
\end{proof}

For integer arguments $d,t,v$, define
\begin{equation}
 \begin{aligned}
 \mathcal R(d,t,v)={}&
 \bigl((t+dv)^3-d^2(t+dv)-d^3\bigr)\\
 &\hspace{1.5em}\cdot\bigl((t+dv)^2+23d^2\bigr).
 \end{aligned}
 \label{i17:homogeneous}
\end{equation}
This is homogeneous of degree five in $d$ and $t$, with $v$ treated
as a parameter.

\begin{lemma}[Combining divisibility with nonvanishing]
\label{i17:merge}
For integers $d,t,z$ with $d\ne0$,
\begin{equation}
 (d\mid t\ \land\ z\ne0)
 \quad\Longleftrightarrow\quad
 \exists v\in\Z:\ zd^5\mid\mathcal R(d,t,v).
 \label{i17:merge-equivalence}
\end{equation}
Moreover, $\mathcal R(d,t,v)\ne0$ for every integer $v$.
\end{lemma}
\begin{proof}
The identity
$\mathcal R(d,t,v)=d^5\Pi(t/d+v)$ and
Lemma~\ref{i17:intersective} prove nonvanishing.
If $t=dq_0$ and $z\ne0$, choose $v$ such that
$z\mid\Pi(q_0+v)$, again using Lemma~\ref{i17:intersective}.
Then $zd^5\mid\mathcal R(d,t,v)$.

Conversely, if $z=0$, the asserted divisibility would mean
$\mathcal R(d,t,v)=0$, which is impossible.
If $d\nmid t$, choose a prime $p$ with
$\alpha=v_p(d)>\beta=v_p(t)$; here $t\ne0$.
For $r=t+dv$ we have $v_p(r)=\beta$.
In $r^3-d^2r-d^3$, the first term has the unique least valuation
$3\beta$. In $r^2+23d^2$, the first term has the unique least
valuation $2\beta$. Hence
\[
 v_p\bigl(\mathcal R(d,t,v)\bigr)
 =5\beta<5\alpha=v_p(d^5),
\]
contrary to the asserted divisibility.
\end{proof}

The general intersectivity mechanism is classical. The argument above
supplies the degree-five instance needed here without making a novelty
claim for $\Pi$ itself. Sun's recent construction uses a related
degree-six divisibility/nonvanishing device \cite[Section 5]{sun2026}.

\subsection{The actual polynomial and its square finalizer}

The quantified coordinates are exactly
\[
 \boldsymbol\xi=
 (\varepsilon,g_0,h_0,s_0,w_0,\phi,i_0,j_0,
   z_1,\ldots,z_7,v,u)\in\Z^{17}.
\]
There are eight base coordinates, seven Pell witnesses, one residue
coordinate $v$, and one final norm coordinate $u$.
Using \eqref{i17:stack}, define
\begin{align*}
 C_\ast&=\Gamma_7,& w_1&=1+C_\ast^2,\\
 w_2&=w_1(P^2+K^2+2),& Z&=(i_0+1)z_1\cdots z_7,\\
 S&=ZD_0^5,& T&=\mathcal R(D_0,T_0,v),\\
 X&=Su-T.
\end{align*}
Then put
\begin{align*}
 A_0&=S^2C_\ast,&
 A_1&=S^2w_1^2a_0,&
 A_2&=S^2w_2^2a_1,\\
 V&=X^2+A_2-A_0-A_1,&
 W&=V^2+4X^2A_2-4A_0A_1.
\end{align*}
The required fixed polynomial is
\begin{equation}
 \boxed{P_{17}=W^2-16V^2X^2A_2.}
 \label{i17:polynomial}
\end{equation}
Every intermediate symbol is a polynomial abbreviation, so none adds
a quantified coordinate. The circuit in
\path{signed17/generated/signed17.json} expands these abbreviations into
a finite division-free directed acyclic graph.

\begin{lemma}[The three-root finalizer]
\label{i17:finalizer}
With the other sixteen coordinates fixed, the assertion
\[
 \exists u\in\Z:\ P_{17}(\boldsymbol\xi)=0
\]
is equivalent to
\begin{equation}
 S\ne0,\qquad S\mid T,\qquad
 C_\ast,a_0,a_1\in\operatorname{Sq}(\Z).
 \label{i17:finalizer-conditions}
\end{equation}
In particular, these conditions are necessary for every zero; conversely,
given these conditions, an integer $u$ can be chosen.
\end{lemma}
\begin{proof}
Lemmas~\ref{i17:denominator} and~\ref{i17:merge} give $T\ne0$.
If $S=0$, then $A_0=A_1=A_2=0$, $X=-T$, and
$P_{17}=T^8\ne0$. Thus a zero must have $S\ne0$.

Choose complex roots $r_0^2=C_\ast$, $r_1^2=a_0$, and $r_2^2=a_1$.
Direct multiplication gives the polynomial identity
\begin{equation}
 P_{17}=\prod_{\epsilon\in\{-1,1\}^3}
 \bigl[X-S(\epsilon_0r_0+\epsilon_1w_1r_1+\epsilon_2w_2r_2)\bigr].
 \label{i17:norm-identity}
\end{equation}
One can verify it first by pairing the signs of $r_0,r_1$, obtaining
$(Y^2-A_0-A_1)^2-4A_0A_1$, and then substituting
$Y=X\pm Sw_2r_2$ and multiplying the two results.
Therefore a zero makes one of the signed sums in
\eqref{i17:norm-identity} equal to the rational number $u-T/S$.

For every integer $C_\ast$,
$\sqrt{|C_\ast|}\le C_\ast^2$, including $C_\ast=0$.
For all real $P,K$,
\begin{equation}
 \sqrt{\left|(P^2-1)K^2+1\right|}\le P^2+K^2+1.
 \label{i17:root-bound}
\end{equation}
Indeed,
$|(P^2-1)K^2+1|\le(P^2+1)K^2+1$, and subtraction from
$(P^2+K^2+1)^2$ leaves
$P^4+K^4+P^2K^2+2P^2+K^2\ge0$.
Consequently
\begin{equation}
 w_1>\sqrt{|C_\ast|},\qquad
 w_2>\sqrt{|C_\ast|}+w_1\sqrt{|a_0|}.
 \label{i17:separation}
\end{equation}

Consider two sign patterns that differ on at least one nonzero root,
and take the largest index at which they differ on a nonzero root.
The absolute value of that root is at least $1$, since its nonzero
square is an integer. By \eqref{i17:separation}, its weighted
contribution to the difference strictly exceeds the sum of the absolute
values of all earlier contributions. The two sums are therefore
distinct. This triangle-inequality argument applies equally to negative
radicands, whose chosen roots are complex.

The field generated by the three roots is a finite Galois extension of
$\mathbb Q$. Every automorphism changes root signs. A rational encoded
sum is fixed by every automorphism, so the separation just proved
forces each root in that sum to be fixed individually.
Every root is therefore rational. A rational number whose square is an
integer is an integer; hence $C_\ast,a_0,a_1$ are nonnegative integer
squares. Zero radicands cause no difficulty, because their roots are
already zero and their signs can be ignored.
The encoded sum is now an integer, so $T/S$ is an integer as well.
This proves necessity of \eqref{i17:finalizer-conditions}.

Conversely, choose integer roots of the three squares and set
\[
 u=T/S+r_0+w_1r_1+w_2r_2\in\Z.
\]
The corresponding factor in \eqref{i17:norm-identity} vanishes.
\end{proof}

\subsection{Equivalence to the fixed consistency statement}

\begin{theorem}[The seventeen-integer equation]
\label{i17:main}
With the coefficients fixed by the recipe above,
\[
 \mathrm{ZFC}\vdash
 \left[
 \left(\nexists\boldsymbol\xi\in\Z^{17}:P_{17}(\boldsymbol\xi)=0\right)
 \longleftrightarrow\Con(\mathrm{ZFC})
 \right].
\]
\end{theorem}
\begin{proof}
A zero gives \eqref{i17:finalizer-conditions}.
Since $D_0\ne0$, the condition $S\ne0$ gives $i\ne0$ and
$z_j\ne0$ for every $j$.
Lemma~\ref{i17:merge} gives $D_0\mid T_0$.
Lemma~\ref{i17:pell-stack} gives the six nonnegative base coordinates
and $m>0$. Lemma~\ref{i17:h-guard} then gives $h_0\ge0$.
Finally, Lemma~\ref{i17:signed-i} permits replacement of $i$ by $|i|$,
yielding $i_0\ge0$ without changing the ratio conditions.
Thus the original eight-natural-variable predicate
\eqref{i17:JR} has a solution. Its imported equivalence supplies a
finite ZFC refutation.

Conversely, a finite ZFC refutation supplies a source solution of
\eqref{i17:JR}, for which $i>0$.
Lemma~\ref{i17:pell-stack} supplies nonzero $z_1,\ldots,z_7$ and
a square $C_\ast$. Lemma~\ref{i17:merge} supplies $v$ with $S\mid T$;
also $S\ne0$. Lemma~\ref{i17:finalizer} supplies $u$, giving an
integer zero of $P_{17}$.

All new steps are elementary number theory and finite Galois theory
formalizable in ZFC. Negating the proved equivalence between existence
of a zero and existence of a refutation gives the assertion.
\end{proof}

This theorem proves the required equivalence. It does not prove either
consistency or inconsistency of ZFC.

\subsection{Exact complexity and reproducible evidence}

The delivered circuit has exactly $20\,607$ binary operations:
$672$ additions, $26$ subtractions, and $19\,909$ multiplications.
Each multiplication by a numeral and each squaring is charged.
As in package 18, construction of the fixed coefficients is outside
the arithmetic evaluation count. The $348$ named fixed leaves are all
bound by the included coefficient recipe.

The exact total degree is
\begin{equation}
 \deg P_{17}=424296\,5^{33}+53824
 =49\,394\,555\,389\,881\,134\,033\,203\,178\,824.
 \label{i17:degree}
\end{equation}
Thus this variable-count improvement has an enormous degree; the
moderate-degree $26$-integer forms retain a substantial advantage on
degree and coefficient complexity.

For a degree audit, recall $L=5^{33}$. The source degrees are
\begin{align*}
 \deg m&=2221L+11,&\deg a_0&=3800L+14,\\
 \deg a_1&=24640L+68,&\deg D_0&=6875L+20.
\end{align*}
The first six guard products have degree recurrence
$c_j=3c_{j-1}+3$, with $c_0=0$, so
$\deg\Gamma_6=1092$. The remaining degrees are
\begin{center}
\begin{tabular}{@{}lr@{}}
\toprule
Expression & Exact degree\\
\midrule
$C_\ast$ & $2221L+3289$\\
$w_1$ & $4442L+6578$\\
$w_2$ & $6342L+6586$\\
$S$ & $34375L+108$\\
$T$ & $34375L+105$\\
$X$ & $34375L+109$\\
$A_2$ & $106074L+13456$\\
\bottomrule
\end{tabular}
\end{center}
The degree of $A_2$ strictly exceeds the degrees of $X^2,A_0,A_1$.
Consequently the leading homogeneous parts of $V$, $W$, and $P_{17}$
are respectively those of $A_2$, $A_2^2$, and $A_2^4$.
The source and stack leading coefficients are nonzero, giving
\eqref{i17:degree}. The generator also propagates structural degree
bounds and a diagonal leading coefficient modulo $1\,000\,003$.
The output coefficient is $40\,718$, is nonzero, and is certified
independent of every fixed coefficient port.

The deterministic finite diagnostics include $24$ modular comparisons
using the actual original T4 coefficients, $150$ independent evaluations
of the eight-factor norm identity, roots of $\Pi$ modulo every modulus
from $1$ through $1200$, $10\,098$ consecutive-square-gap checks, and
the signed-$i$ symmetry. These checks supplement the universal proofs
above. Run
\begin{lstlisting}
python3 -B signed17/build_signed17.py
\end{lstlisting}
from the archive root. The standard-library generator writes
\path{signed17/generated/signed17.json} and
\path{signed17/generated/signed17_checks.json}.
The latter is a diagnostic receipt, not a proof-assistant certificate.

\subsection{Provenance and remaining opportunities}

The unchanged source expressions and coefficient modules come from
package 18 at revision \pin. That package records the earlier source
acquisition revision
\texttt{f0ae4e0991ca44a7bb4e5671d42e33d775e48035}.
All nineteen imported files are authenticated byte for byte against
their pinned Git blob identifiers, with SHA-256 hashes recorded in
\path{signed17/SOURCE_PINS.json}. The literal quartic and ordered tile
list also match the hashes in the upstream coefficient binding.
The generator checks every source pin before rebuilding the polynomial.
The archive preserves the upstream license and attribution files.

Several questions remain open for further simplification. Recovering
the $j_0$ or $\phi$ nonnegativity condition from the other source
conditions would remove another guard witness. A complete fixed
instantiation of the published ten-integer construction would improve
the variable count further. Optimizing the digit-polynomial circuit
could reduce its roughly $20\,500$ inherited operations. Finally, a
rebinding to the smaller $943$-tile tuple would require proof of the
complete quartic/compiler interface for that tuple; individual
coefficient ports cannot be spliced between the two bindings.
The present section makes no exhaustive novelty or absolute-minimality
claim for its general devices.
